\scanpage{58}
\mspage{II. 6}
% top insertion, continued in slanted left margin
\inserted{Il existe donc un \unsure{ss-gro|ss-groupe} ouvert $\pi'\hookrightarrow\pi_1(X,\unsure{\xi|\zeta})$ sur lequel la repr. est unipotente. Mais $\pi'$ définit un rev. étale connexe $X'$ de $X$, \deleted{\unsure{tel}} $\xi'$ pt au-dessus tel que la représentation naturelle de $\pi_1(X',\xi')=\pi$ dans $T_\ell(A_{X'}(\xi'))$ soit unipotente. On peut donc, pour la démonstration, supposer déjà que $\pi_1(X,\xi)$ opère de façon unipotente.}
% slanted: continuation of the insertion in the left margin

% lines 3-8: crossed-out argument
\deleted{Supposons maintenant que la représentation envisagée de $\pi_1(X,\xi)$ soit aussi semi-simple, donc, étant unipotente, elle est triviale sur un ss-groupe \deleted{distingué} \inserted{ouvert} $U$ de $\pi_1(X,\xi)$. Ce sous-groupe définit un revêtement \inserted{connexe $\xi$-pointé} $X'$ de $X$, et par construction la représentation de $\pi_1(X',\xi)$ dans $T_\ell(A_{\xi'})$ est triviale.}

% lines 9-13
\deleted{\unsure{Cela}} Il suffit alors d'appliquer \inserted{l'adaptation \unsure{suivante|immédiate} d'}un résultat bien connu de \unsure{descente}, \deleted{\unsure{simple}} utilisant Mordell-Weil (cf Inventiones, Un th. sur les hom. de schémas abéliens, prop. 4.4) : \deleted{\unsure{on en déduit aussitôt que}}

\uline{Lemme} 1.5.1. \inserted{Si l'action de $\pi_1(X,\xi)$ sur $T_\ell(A(\xi))$ est \deleted{\unsure{ess.}} unipotente,} \deleted{\unsure{Si}} Alors cette action est triviale, et $\exists$ schéma abélien $B$ sur $k$, et \inserted{une isogénie radicielle $B_X\longrightarrow A$}.
% slanted: the isogeny insertion is in the left margin

% lines 14-18
\uline{Corollaire.} 1.6. \deleted{Sous les hyp.} Si pour \uline{un} $\ell\ne\operatorname{car}k$, l'action de $\pi_1(X,\xi)$ sur $T_\ell(A_\xi)$ est \deleted{unipotente \unsure{et}} \inserted{ess\textsuperscript{t} résoluble} \deleted{semi-simple} ($\Longleftrightarrow$ \deleted{\unsure{quasi}} essentiellement triviale) il en est de même pour tout autre $\ell'\ne\operatorname{car}k$.

% lines 19-26
\uline{Corollaire} 1.7. Supposons \deleted{que} \inserted{$\operatorname{car}k=0$, et} \deleted{\unsure{que}} l'action de $\pi_1(X,\xi)$ sur $T_\ell(A_\xi)$ soit \unsure{abélienne|unipotente} et semi-simple, et supposons qu'il existe un \inserted{entier $n\geq3$}, \deleted{\unsure{premier} $\ell'\ne\operatorname{car}k$,} tel que la représentation \deleted{de} induite \inserted{de} $\pi_1(X,\xi)$ sur ${}_n A(\xi)$ soit triviale, i.e. que ${}_n A$ soit un revêtement en groupes \uline{trivial} de $X$. Alors \deleted{\unsure{en 1.5 on peut prendre} $X'\longrightarrow X$, (en particulier si $\operatorname{car}k=0$)} $A$ est un schéma abélien \uline{constant} sur $X$.

% lines 27-29
En effet, \deleted{\unsure{alors} $\pi_1(X,\xi)$ opère sur $T_\ell(A_\xi)$ par l'intermédiaire d'un groupe \uline{fini}. Or on} la théorie de la descente galoisienne et la structure discrète de $\operatorname{Aut}_k(B)$ impliquent que
% archive sheet number: 57
